\chapter{ضرب الأعداد السالبة}\label{ch:g8:negprod}

جمعت السنة السابعة الأعداد النسبية وطرحتها
(\cref{ch:g7:negatives})؛ ويبقى ضربها وقسمتها. ويهيمن
سؤال واحد على الفصل: لماذا يكون «ناقص في ناقص»
زائدًا؟ ونحن نعطي سببًا، لا قاعدة فحسب.

\section{قواعد الإشارة}

\begin{theorem}[إشارة الجداء]\label{thm:g8:negprod:rules}
\omterm{def:g2:mult:def}{جداء} عددين نسبيين مسافته إلى \omterm{def:g1:counting:zero}{الصفر} تساوي
\omterm{def:g2:mult:def}{جداء} المسافتين، وإشارته معطاة بما يلي:
\begin{center}
\renewcommand{\arraystretch}{1.3}
\begin{tabular}{c|cc}
$\times$ & $+$ & $-$ \\
\hline
$+$ & $+$ & $-$ \\
$-$ & $-$ & $+$ \\
\end{tabular}
\end{center}
\emph{الإشارتان نفسهما: \omterm{def:g2:mult:def}{جداء} موجب. وإشارتان متعاكستان: \omterm{def:g2:mult:def}{جداء} سالب.}
والجدول نفسه يحكم \omterm{def:g3:division:remainder}{خوارج القسمة}.
\end{theorem}

\begin{proof}[لماذا $(-1) \times (-1) = +1$]
أولًا، $3 \times (-4)$ يعني $(-4) + (-4) + (-4) = -12$: فموجب
في سالب سالب. والآن راقب النمط:
\[
3 \times (-4) = -12, \quad
2 \times (-4) = -8, \quad
1 \times (-4) = -4, \quad
0 \times (-4) = 0 .
\]
في كل خطوة، ينقص العامل الأول بمقدار $1$ و\emph{ترتفع} النتيجة
بمقدار $4$. وبمواصلة النمط خطوة أخرى:
\[
(-1) \times (-4) = 0 + 4 = +4 .
\]
وأيّ اِختيار آخر يكسر اِنتظام الحساب
(التوزيعية). إذن ناقص في ناقص لا بدّ أن يكون زائدًا.
\end{proof}

\begin{omfigure}
\begin{tikzpicture}
\begin{axis}[omaxis,
  width=7.8cm, height=5.4cm,
  xmin=-2.6, xmax=3.6, ymin=-14, ymax=10,
  xlabel={$k$}, ylabel={$k \times (-4)$},
  xtick={-2,-1,1,2,3}, ytick={-12,-8,-4,4,8}]
  \addplot[omDef, thick, domain=-2.3:3.3] {-4*x};
  \addplot[omDef, only marks, mark=*, mark size=1.6pt] coordinates
    {(3,-12) (2,-8) (1,-4) (0,0) (-1,4) (-2,8)};
\end{axis}
\end{tikzpicture}

{\small حجّة النمط في صورة: فالنقط
$(k,\ k \times (-4))$ تصطفّ، ومواصلة \omterm{def:g6:lines:objects}{المستقيم} بعد $k = 0$
تفرض $(-1)\times(-4) = +4$ و $(-2)\times(-4) = +8$.}
\end{omfigure}

\begin{example}\label{ex:g8:negprod:products}
\begin{align*}
(-7) \times (+6) &= -42 && \text{(إشارتان متعاكستان)} \\
(-5) \times (-8) &= +40 && \text{(الإشارتان نفسهما)} \\
(+2.5) \times (-4) &= -10 \\
(-36) \div (-9) &= +4 \\
(+15) \div (-2) &= -7.5 .
\end{align*}
\end{example}

\section{جداءات عوامل عدّة}

\begin{proposition}[إشارة جداء طويل]\label{prop:g8:negprod:many}
إشارة \omterm{def:g2:mult:def}{جداء} عوامل عدّة غير معدومة تتوقف على عدد
العوامل السالبة وحده:
\begin{itemize}
  \item عدد \emph{زوجي} من العوامل السالبة يعطي \omterm{def:g2:mult:def}{جداء}
        موجبًا؛
  \item وعدد \emph{فردي} يعطي \omterm{def:g2:mult:def}{جداء} سالبًا.
\end{itemize}
\end{proposition}

\begin{proof}
تتزاوج العوامل السالبة، وكل زوج يسهم بإشارة موجبة
(\cref{thm:g8:negprod:rules})؛ \omterm{def:g3:numbers:evenodd}{والعدد الفردي} يترك عاملًا سالبًا
بلا زوج، فيجعل \omterm{def:g2:mult:def}{الجداء} كله سالبًا.
\end{proof}

\begin{example}\label{ex:g8:negprod:many}
$(-2) \times (+3) \times (-5) \times (-1)$: ثلاثة عوامل سالبة
(\omterm{def:g3:numbers:evenodd}{عدد فردي})، إذن \omterm{def:g2:mult:def}{الجداء} سالب؛ والمسافات: $2 \times 3 \times 5
\times 1 = 30$. والنتيجة: $-30$. فحدّد الإشارة \emph{أولًا}، ثم
اِضرب المسافات --- مهمتان سهلتان بدل واحدة معرّضة
للخطأ.
\end{example}

\begin{example}[قوى الأعداد السالبة]\label{ex:g8:negprod:powers}
$(-2)^2 = (-2) \times (-2) = +4$، و
$(-2)^3 = (-2)^2 \times (-2) = -8$: فالأُسُس الزوجية تعطي قيمًا
موجبة، والأُسُس الفردية تحفظ الإشارة. واِنتبه للترميز:
$(-3)^2 = 9$ لكن $-3^2 = -(3 \times 3) = -9$.
\end{example}

\section{الحساب بالعمليات الأربع}

\begin{method}[حساب آمن بالأعداد النسبية]\label{met:g8:negprod:compute}
\begin{enumerate}
  \item اِحترم أولويات \cref{ch:g7:priorities}: الأقواس،
        ثم $\times$ و $\div$، ثم $+$ و $-$؛
  \item وعند كل ضرب أو قسمة، جِد \emph{الإشارة
        أولًا}، ثم المسافة؛
  \item وأعِد الكتابة خطوة واحدة في كل سطر.
\end{enumerate}
\end{method}

\begin{example}\label{ex:g8:negprod:compute}
\begin{align*}
5 - 3 \times (-4)
&= 5 - (-12) && \text{(الجداء أولًا: إشارتان متعاكستان)}\\
&= 5 + 12 = 17 .
\end{align*}
\begin{align*}
\frac{(-6) \times (-10)}{-4}
&= \frac{+60}{-4} && \text{(البسط: الإشارتان نفسهما)}\\
&= -15 && \text{(خارج القسمة: إشارتان متعاكستان).}
\end{align*}
\end{example}

\begin{example}[التعويض بقيم سالبة]\label{ex:g8:negprod:substitute}
اِحسب قيمة $E = 3x^2 - 2x$ من أجل $x = -5$، مع أقواس حول القيمة:
\[
E = 3 \times (-5)^2 - 2 \times (-5)
= 3 \times 25 - (-10)
= 75 + 10 = 85 .
\]
والأقواس في $(-5)^2$ ضرورية: فبدونها ينطبق التربيع
على $5$ وحده.
\end{example}

\section{تمارين}

\begin{exercise}[$\star$]\label{exo:g8:negprod:1}
اِحسب:
\[
(-8) \times (+7), \qquad
(-9) \times (-6), \qquad
(+12) \times (-0.5), \qquad
(-1) \times (-1) \times (-1).
\]
\end{exercise}

\begin{exercise}[$\star$]\label{exo:g8:negprod:2}
اِحسب:
\[
(-63) \div (+9), \qquad
(-8) \div (-16), \qquad
\frac{45}{-5}, \qquad
\frac{-7.2}{-8} .
\]
\end{exercise}

\begin{exercise}[$\star$]\label{exo:g8:negprod:3}
أعطِ \emph{إشارة} كل \omterm{def:g2:mult:def}{جداء} فقط، دون حسابه:
\[
(-3) \times (-7) \times (+2) \times (-5); \qquad
(-1)^{10}; \qquad
(-2)^{7}; \qquad
(-4) \times 0 \times (-6).
\]
\end{exercise}

\begin{exercise}[$\star$]\label{exo:g8:negprod:4}
اِحسب:
\[
(-2)^4, \qquad
-2^4, \qquad
(-3)^3, \qquad
(-1)^{2026} .
\]
\end{exercise}

\begin{exercise}[$\star$]\label{exo:g8:negprod:5}
اِحسب خطوة خطوة، مع اِحترام الأولويات:
\[
7 + 2 \times (-6), \qquad
(-4) \times (5 - 9), \qquad
18 \div (-3) - (-2) .
\]
\end{exercise}

\begin{exercise}[$\star$]\label{exo:g8:negprod:6}
اِحسب:
\[
\frac{(-5) \times (+8)}{-4}, \qquad
\frac{(-3) \times (-6)}{(-2) \times (+9)} .
\]
\end{exercise}

\begin{exercise}[$\star$]\label{exo:g8:negprod:7}
اِحسب القيمة من أجل $x = -3$:
\[
5x, \qquad x^2, \qquad -x^2, \qquad 2x^2 + 4x, \qquad (2x)^2 .
\]
\end{exercise}

\begin{exercise}[$\star\star$]\label{exo:g8:negprod:8}
أكمل كل تساوٍ:
\[
(-6) \times \;?\; = 42, \qquad
\;?\; \div (-4) = -2.5, \qquad
(-5) \times \;?\; = -1 .
\]
\end{exercise}

\begin{exercise}[$\star\star$]\label{exo:g8:negprod:9}
صواب أم خطأ؟ بَرِّر أو أعطِ مثالًا مضادًّا.
\begin{enumerate}
  \item \omterm{def:g2:mult:def}{جداء} عددين نسبيين لا يقلّ أبدًا عن
        مجموعهما.
  \item مربع \omterm{def:g7:negatives:def}{عدد نسبي} ليس سالبًا أبدًا.
  \item إذا كان \omterm{def:g2:mult:def}{جداء} ثلاثة عوامل موجبًا، فالعوامل الثلاثة
        كلها موجبة.
\end{enumerate}
\end{exercise}

\begin{exercise}[$\star\star$]\label{exo:g8:negprod:10}
كل جواب خاطئ في اِختبار يساوي $-3$ نقاط، وكل جواب صحيح
$+5$. أجابت زينة عن $20$ سؤالًا فحصلت على $36$ نقطة. فكم
جوابًا كان صحيحًا؟ (جرّب قيمًا، أو أقِم الحساب
$5r - 3(20 - r) = 36$.)
\end{exercise}

\begin{exercise}[$\star\star\star$]\label{exo:g8:negprod:11}
باستعمال التوزيعية (كما في برهان
\cref{thm:g8:negprod:rules})، اُنشر وبَرِّر كل خطوة من
\[
0 = (-1) \times 0 = (-1) \times \bigl(1 + (-1)\bigr)
= (-1) \times 1 + (-1) \times (-1),
\]
واِستنتج أن $(-1) \times (-1)$ لا بدّ أن يساوي $+1$.
\end{exercise}

\section{مسألة: لماذا يكون ناقص في ناقص زائدًا}

\begin{problem}[{مسألة نهاية الأسبوع --- قواعد الإشارة مستنتجة من
التوزيعية، والمجموع المتناوب $1 - 2 + 3 - 4 + \dots$}]
\label{pb:g8:negprod:1}
واصل برهان \cref{thm:g8:negprod:rules} نمطًا
واِدّعى أن «أيّ اِختيار آخر يكسر التوزيعية». وهذه
المسألة تجعل الادّعاء مضبوطًا: فاِنطلاقًا من التوزيعية وحدها
(\cref{thm:g7:priorities:distributivity})، \emph{ستبرهن} على
قواعد الإشارة --- بلا صورة ولا نمط، كما يفعل الجبر ---
ثم تشغّلها على قوى العدد $-1$ وعلى \omterm{def:g1:addition:def}{مجموع} مشهور ذي
إشارات متناوبة. وفي كل ما يلي، $a$ و $b$ عددان نسبيان؛ ونسلّم
بقواعد الجمع في \cref{ch:g7:negatives}، وبأن
$1 \times a = a$، وبأن \omterm{def:g2:mult:def}{الجداء} يمكن حسابه بأيّ ترتيب
(فتنطبق التوزيعية على أيّ من طرفَي \omterm{def:g2:mult:def}{الجداء}).

\medskip
\noindent\textbf{الجزء الأول --- قواعد الإشارة مبرهنات.}
\begin{enumerate}
  \item بَرهِن على أن $0 \times a = 0$ من أجل كل \omterm{def:g7:negatives:def}{عدد نسبي} $a$.
        (اُكتب $0 = 0 + 0$، واُنشر $(0 + 0) \times a$، واِسأل أيّ
        عدد، إذا جُمع إلى $0 \times a$، يعيد $0 \times a$.)
  \item اُنشر $\bigl(1 + (-1)\bigr) \times a$ لتبيّن أن
        $(-1) \times a + a = 0$، واِستنتج أن
        \[
        (-1) \times a = -a :
        \]
        فالضرب في $-1$ يعطي \omterm{def:g7:negatives:def}{المعاكس}. وقارن مع
        \cref{exo:g8:negprod:11}، وهو الحالة $a = -1$.
  \item اِستنتج أن $(-a) \times b = -(a \times b)$: فإشارة ناقص
        واحدة تخرج من \omterm{def:g2:mult:def}{الجداء} دون تغيير.
  \item واِستنتج أن $(-a) \times (-b) = a \times b$: فإشارتا ناقص
        تتلاشيان.
  \item كل \omterm{def:g7:negatives:def}{عدد نسبي} هو مسافته إلى \omterm{def:g1:counting:zero}{الصفر} مع إشارة
        أمامها: $a = +d$ أو $a = -d$. اِشرح كيف يبرهن السؤالان 3 و~4
        على حالات جدول الإشارة الأربع كلها في
        \cref{thm:g8:negprod:rules}، بما فيها المسافات.
\end{enumerate}

\medskip
\noindent\textbf{الجزء الثاني --- قوى العدد $-1$.}
\begin{enumerate}[resume]
  \item اِحسب $(-1)^2$ و $(-1)^3$ و $(-1)^4$ و $(-1)^5$، ثم أعطِ،
        مع التبرير من \cref{prop:g8:negprod:many}،
        قيمة $(-1)^n$ من أجل كل عدد صحيح $n \geq 1$.
  \item أعطِ (دون حساب أيّ مسافة) إشارة
        \[
        (-1) \times (-2) \times (-3) \times \dots \times (-10),
        \]
        واُكتب مسافة هذا \omterm{def:g2:mult:def}{الجداء} جداءَ أعداد
        صحيحة، دون إنجازه.
  \item بيّن أن $(-a)^2 = a^2$ وأن $(-a)^3 = -a^3$. فمن أجل أيّ
        أُسُس تصمد إشارة الناقص؟
  \item اِستعمل قواعد الإشارة لتبيّن أن مربع \omterm{def:g7:negatives:def}{عدد
        نسبي} ليس سالبًا أبدًا، واِستنتج أنه لا يوجد \omterm{def:g7:negatives:def}{عدد نسبي}
        $x$ يحقّق $x^2 = -4$.
  \item تدّعي زينة: «$(-1)^{2026} + (-1)^{2027} = 0$، وعمومًا
        تتلاشى قوّتان متتاليتان للعدد $-1$ دائمًا.» فهل هي
        على صواب؟ بَرِّر.
\end{enumerate}

\medskip
\noindent\textbf{الجزء الثالث --- \omterm{def:g1:addition:def}{المجموع} المتناوب.}
بعد تأمين قواعد الإشارة، يمكننا حساب مجاميع تخلط
الإشارتين. ومن أجل عدد صحيح $n \geq 1$، ليكن
\[
A_n = 1 - 2 + 3 - 4 + \dots
\]
\omterm{def:g1:addition:def}{مجموع} الأعداد الصحيحة من $1$ إلى $n$ بإشارات
متناوبة: فالحدّ الأخير $+n$ حين يكون $n$ فرديًّا، و $-n$ حين يكون $n$
زوجيًّا. ومنه $A_1 = 1$ و $A_2 = 1 - 2 = -1$.
\begin{enumerate}[resume]
  \item اِحسب $A_3$ و $A_4$ و $A_5$ و $A_6$ و $A_7$. فماذا
        تخمّن؟
  \item اِفترض $n$ زوجيًّا. جمّع الحدود مثنى مثنى،
        $(1 - 2) + (3 - 4) + \dots$، وعُدّ الأزواج، وبَرهِن على أن
        $A_n = -\frac{n}{2}$.
  \item واِفترض $n$ فرديًّا. اِشرح لماذا $A_n = A_{n-1} + n$،
        واِستنتج من السؤال السابق أن
        $A_n = \frac{n + 1}{2}$.
  \item اِحسب $A_{100}$ و $A_{101}$ و $A_{2026}$.
  \item حسبت عالية $A_{101} - A_{100} = 51 - (-50) = 101$ ففوجئت
        بأن الناتج $101$ بالضبط. فهل كان عليها أن تُفاجأ؟
        اِشرح في جملة واحدة، واُذكر قيمة
        $A_{2027} - A_{2026}$ دون أيّ حساب آخر.
\end{enumerate}
\end{problem}
